\scanpage{163}
\mspage{79}
% archive sheet number: 162

4.4. \uline{Étude \unsure{de|des}} $\operatorname{Ind}_\pi(C)$ (pour $C$ \unsure{petite} catégorie stable par $\varinjlim$ de cardinal $\leq\pi$).

4.4.1. Finalement, la démonstration du lemme 4.3.7 n'est pas tellement évidente que je l'\unsure{anti}cipais, et il me faut lui consacrer des préliminaires techniques, que je \unsure{réunirai|vais réunir} ici. Dans tte la suite, $C$ représente une \inserted{\unsure{$U$|$\mathcal U$}-}catégorie équivalente à une petite catégorie.

\uline{Proposition 4.4.2.} Soit $F\in C^\wedge$ un foncteur ind-représentable, i.e tel que $C_{/F}$ soit filtrante. Soit d'autre part $i\longmapsto x_i$ un foncteur $\unsure{\varphi|q}\,I\longrightarrow C$, où $I$ est une petite catégorie filtrante, et \deleted{considérons} soit

\[
\alpha:x=\varinjlim_I{}_{\!(C^\wedge)}x_i\longrightarrow \deleted{\unsure{E|F}}F
\tag*{(4.4.2.1)}
\]

% slanted
\margin{La donnée de $\alpha$ équivaut aussi à celle d'une flèche de ind-objets de $(x_i)_{i\in I}\longrightarrow F$, $F$ étant considéré comme un Ind-objet \deleted{\unsure{constant}} de $C$ \unsure{indexé par} $C/F$.}

un système projectif de flèches $x_i\unsure{-|\longrightarrow} F$ (NB, la limite inductive $x$ est prise dans $C^\wedge$, et est calculée argument par \uline{argument}), \inserted{ce qui équivaut à} \deleted{de sorte que $\unsure{\varphi|q}\,I\longrightarrow C$ \unsure{se}} une factorisation $\varphi:I\longrightarrow C_{/F}$ de $I\longrightarrow C$ en

\diagram{\begin{tikzcd}[ampersand replacement=\&]
I \arrow[r,"\varphi_\alpha"] \& C_{/F} \arrow[r,hook,"\mathrm{can}.", "\unsure{\psi|+|r}"'] \& C.
\end{tikzcd}}

Ceci dit, pour que $\alpha$ (4.4.\unsure{2|?}.1) soit un isomorphisme, dans $\operatorname{Ind}'(C)\subset C^\wedge$ (ss-catégorie pleine des foncteurs ind-représentables sur $C$), il f. et s. que $\varphi_\alpha$ soit \uline{cofinal} (ou ce qui revient au $=$, $I$ étant filtrant) qu'il satisfasse les deux conditions habituelles

\[
\begin{aligned}
\mathrm{F1})\quad &\forall\ \unsure{\zeta|\xi}\in\operatorname{Ob} C_{/F},\ \exists i\in\operatorname{Ob} I\text{ et une flèche }\zeta\longrightarrow\varphi_\alpha(i)\\
\mathrm{F2})\quad &\forall\ \deleted{\unsure{z}}\zeta\in\operatorname{Ob} C_{/F},\ \text{et }i\in\operatorname{Ob} I,\text{ et tte double flèche}
\end{aligned}
\]

